\section{Results}

\subsection{Positional Distribution of Uncertainty (RQ1)}

Uncertainty expression follows a clear inverted-U trajectory across reasoning chain positions (Figure~\ref{fig:positional}, Table~\ref{tab:deciles}). The rate is lowest at the beginning of the chain (2.4\% in decile 1), rises sharply to a peak at decile 3 (6.9\%), then gradually declines toward the conclusion (4.3\% in decile 10). This pattern is consistent across all six uncertainty categories, though with different magnitudes.

\begin{figure}[t]
  \centering
  \includegraphics[width=\linewidth]{fig1_uncertainty_density.pdf}
  \caption{Uncertainty expression rate across reasoning chain position deciles. The inverted-U pattern peaks at decile 3 and declines toward conclusions.}
  \label{fig:positional}
\end{figure}

\begin{table}[t]
\centering
\caption{Uncertainty expression rate by position decile in reasoning chains. Rates represent the proportion of sentences containing at least one uncertainty marker. The overall rate peaks in deciles 2--4 and declines toward the conclusion.}
\label{tab:deciles}
\small
\begin{tabular}{rrrrrrrr}
\toprule
Decile & $N$ & Overall & Epist. & Evid. & Prob. & Modal & Approx. \\
\midrule
1 & 25{,}864 & 2.4 & 0.8 & 0.12 & 0.8 & 0.8 & 0.10 \\
2 & 19{,}178 & 6.0 & 2.4 & 0.10 & 1.4 & 2.1 & 0.21 \\
\textbf{3} & 19{,}950 & \textbf{6.9} & 2.8 & 0.14 & 1.4 & 2.6 & 0.21 \\
4 & 19{,}443 & 6.4 & 2.6 & 0.10 & 1.1 & 2.6 & 0.15 \\
5 & 18{,}536 & 6.3 & 2.4 & 0.04 & 1.0 & 2.8 & 0.27 \\
6 & 21{,}046 & 5.8 & 2.3 & 0.18 & 0.9 & 2.6 & 0.25 \\
7 & 20{,}176 & 5.9 & 2.1 & 0.10 & 0.9 & 2.7 & 0.35 \\
8 & 19{,}217 & 5.6 & 2.1 & 0.10 & 0.8 & 2.5 & 0.33 \\
9 & 19{,}911 & 5.8 & 2.7 & 0.11 & 0.7 & 2.3 & 0.18 \\
10 & 24{,}835 & 4.3 & 1.9 & 0.08 & 0.7 & 1.6 & 0.15 \\
\bottomrule
\end{tabular}
\end{table}

Modal hedges and epistemic hedges are the dominant categories throughout, each appearing in approximately 2--3\% of mid-chain sentences. Probability language contributes roughly 1\%, while evidential markers, explicit uncertainty, and approximators are rarer ($< 0.3\%$ each). Figure~\ref{fig:categories} shows the category-specific positional curves.

\begin{figure}[t]
  \centering
  \includegraphics[width=\linewidth]{fig2_category_by_position.pdf}
  \caption{Uncertainty categories across reasoning position. Modal hedges and epistemic hedges dominate; all categories follow the inverted-U pattern.}
  \label{fig:categories}
\end{figure}

The low initial rate (decile 1) likely reflects that reasoning chains often begin with task restatement or setup, which involves less deliberation. The mid-chain peak corresponds to the evaluative phase where the model considers alternatives, weighs evidence, and entertains competing hypotheses. The late-chain decline suggests a commitment phase where the model converges on an answer and reduces hedging.

\subsection{Cross-Model Comparison (RQ2)}

Model families exhibit broadly similar inverted-U positional patterns, but differ substantially in overall uncertainty magnitude (Table~\ref{tab:families}, Figure~\ref{fig:models}).

\begin{figure}[t]
  \centering
  \includegraphics[width=\linewidth]{fig4_model_comparison.pdf}
  \caption{Uncertainty rate by position decile across model families. All families show the inverted-U pattern, but at different magnitudes.}
  \label{fig:models}
\end{figure}

\begin{table}[t]
\centering
\caption{Uncertainty rates by model family. Mean rate is averaged across all position deciles. Families with fewer than 100 total sentences are excluded.}
\label{tab:families}
\small
\begin{tabular}{lrrrr}
\toprule
Model Family & $N$ Sentences & Mean Rate (\%) & Peak Decile & Peak Rate (\%) \\
\midrule
other & 1{,}345 & 22.0 & 5 & 35.0 \\
qwen & 2{,}038 & 13.8 & 6 & 23.0 \\
minimax & 23{,}855 & 12.2 & 7 & 15.3 \\
deepseek & 20{,}097 & 5.7 & 5 & 8.3 \\
claude-sonnet & 1{,}808 & 5.1 & 8 & 9.5 \\
glm & 85{,}735 & 4.8 & 3 & 7.7 \\
gpt & 22{,}966 & 4.2 & 2 & 7.4 \\
claude-opus & 34{,}587 & 3.9 & 7 & 5.0 \\
gemini & 15{,}722 & 2.5 & 4 & 3.8 \\
\bottomrule
\end{tabular}
\end{table}

Some families (e.g., Minimax, Qwen) display notably higher uncertainty rates, exceeding 15\% in peak deciles. In contrast, major families like Claude Opus, GPT, and Gemini show lower overall rates (3--5\% at peak). Despite these magnitude differences, the \emph{shape} of the positional curve is remarkably consistent: all families peak in deciles 2--5 and decline thereafter. The mixed-effects model (Section~\ref{sec:regression}) confirms that model family dummies are not statistically significant predictors of uncertainty when controlling for position and other covariates ($p > .05$ for all families).

\subsection{Confidence Filtering (RQ3)}

\begin{table}[t]
\centering
\caption{Confidence filtering: paired comparison of uncertainty rates in reasoning traces versus final responses (Wilcoxon signed-rank test, one-sided). Significance: {*}$p < .05$, {**}$p < .01$, {***}$p < .001$.}
\label{tab:filtering}
\small
\begin{tabular}{lrrrl}
\toprule
Category & Reasoning (\%) & Response (\%) & $\Delta$ (\%) & Sig. \\
\midrule
\textbf{Overall} & 3.84 & 1.66 & +2.18 & *** \\
Epistemic Hedge & 1.34 & 0.20 & +1.14 & *** \\
Evidential Marker & 0.15 & 0.11 & +0.04 & * \\
Explicit Uncertainty & 0.08 & 0.02 & +0.06 & *** \\
Probability Language & 0.87 & 0.34 & +0.53 & *** \\
Modal Hedge & 1.47 & 0.63 & +0.83 & *** \\
Approximator & 0.21 & 0.42 & -0.21 &  \\
\bottomrule
\end{tabular}
\end{table}

We observe substantial confidence filtering between reasoning traces and final responses (Table~\ref{tab:filtering}). Overall, the uncertainty rate drops from 3.8\% in reasoning to 1.7\% in responses, a reduction of more than half ($W = 4{,}397{,}479$, $p < 10^{-165}$, Wilcoxon signed-rank).

\begin{figure}[t]
  \centering
  \includegraphics[width=\linewidth]{fig3_confidence_filtering.pdf}
  \caption{Uncertainty rates in reasoning traces versus final responses (paired comparison). All categories except approximators show significant reduction.}
  \label{fig:filtering}
\end{figure}

The filtering effect varies by uncertainty category. Epistemic hedges show the strongest suppression, dropping from 1.34\% to 0.20\% ($p < 10^{-189}$). Modal hedges are reduced from 1.47\% to 0.63\% ($p < 10^{-66}$). Probability language drops from 0.87\% to 0.34\% ($p < 10^{-73}$). Explicit uncertainty, already rare, is further suppressed from 0.08\% to 0.02\% ($p < 10^{-10}$). Evidential markers show modest reduction ($p = .014$).

Notably, \emph{approximators} are the exception: they slightly \emph{increase} from 0.21\% in reasoning to 0.42\% in responses ($p = .062$, not significant). This is consistent with the idea that approximators serve a communicative function in final answers (``about 50 results'') rather than reflecting deliberative uncertainty.

\begin{figure}[t]
  \centering
  \includegraphics[width=\linewidth]{fig5_filtering_distribution.pdf}
  \caption{Distribution of per-interaction filtering ratios. Positive values indicate uncertainty reduction from reasoning to response.}
  \label{fig:filtering_dist}
\end{figure}

\subsection{Regression Results}
\label{sec:regression}

The mixed-effects model confirms the positional pattern and identifies additional predictors (Table~\ref{tab:regression}).

\begin{table}[t]
\centering
\caption{Mixed-effects linear model results. Panel (a): positional uncertainty model (DV: uncertainty presence per sentence, 208{,}156 observations across 9{,}923 interactions). Panel (b): filtering ratio model. Model family dummies were included in (a) but are omitted for brevity (all $p > .05$).}
\label{tab:regression}
\small
\begin{tabular}{lrrrr}
\toprule
Predictor & Coef. & Std.Err. & $z$ & $p$ \\
\midrule
\multicolumn{5}{l}{\textit{(a) Positional uncertainty model}} \\
\midrule
Intercept & -0.020 & 0.138 & -0.148 & $0.882$ \\
Position (linear) & 0.137 & 0.006 & 22.343 & $< 0.001$*** \\
Position$^2$ (quadratic) & -0.132 & 0.006 & -22.209 & $< 0.001$*** \\
Log reasoning length & 0.005 & 0.001 & 5.812 & $< 0.001$*** \\
NSFW content & 0.019 & 0.005 & 4.139 & $< 0.001$*** \\
\midrule
\multicolumn{5}{l}{\textit{(b) Filtering ratio model} ($n$ = 9{,}611)} \\
\midrule
Intercept & -0.092 & 0.071 & -1.293 & $0.196$ \\
Log reasoning length & 0.118 & 0.006 & 21.201 & $< 0.001$*** \\
\bottomrule
\end{tabular}
\end{table}

The significant positive coefficient for normalized position ($\beta = 0.137$, $p < .001$) combined with the significant negative quadratic term ($\beta = -0.132$, $p < .001$) formally establishes the inverted-U trajectory. Longer reasoning chains are associated with slightly higher uncertainty ($\beta = 0.005$, $p < .001$), and NSFW content shows elevated uncertainty ($\beta = 0.019$, $p < .001$). Model family effects are uniformly non-significant.

For the filtering ratio model, log reasoning length is the sole significant predictor ($\beta = 0.118$, $p < .001$): longer reasoning chains show more aggressive confidence filtering, suggesting that extended deliberation leads to greater suppression of uncertainty in the final answer.
